\centering
\resizebox{\ifdim\width>\linewidth\linewidth\else\width\fi}{!}{
\begin{tabular}{lcccc}
\toprule
Structural model & Control PoTs & Bracelet PoTs & PoTs saved & PoTs saved, static \\
\midrule
Baseline model (cluster bootstrap) & 109 & 81 & \makecell[c]{25\\(10, 42)} & \makecell[c]{7\\(-5, 23)} \\
Correct classification & 107 & 89 & \makecell[c]{18\\(6, 29)} & \makecell[c]{1\\(-12, 14)} \\
Perceived observability (second order) & 98 & 76 & \makecell[c]{21\\(13, 30)} & \makecell[c]{5\\(0, 12)} \\
Any-signal/no-signal social image weight & 108 & 79 & \makecell[c]{28\\(18, 36)} & \makecell[c]{11\\(-2, 22)} \\
Treatment-specific social image weight & 108 & 83 & \makecell[c]{25\\(-27, 37)} & \makecell[c]{8\\(-42, 26)} \\
Student-$t(5)$ intrinsic-motivation types & 107 & 79 & \makecell[c]{27\\(17, 35)} & \makecell[c]{7\\(1, 16)} \\
Cluster random shock & 107 & 94 & \makecell[c]{12\\(1, 28)} & \makecell[c]{3\\(0, 11)} \\
\bottomrule
\end{tabular}}
\begin{minipage}{0.98\linewidth}
\footnotesize \textit{Notes:} PoTs saved is the draw-level paired difference between the Control and Bracelet optimal allocations at the common welfare target. Parentheses contain 2.5th and 97.5th percentiles. The final column holds social-image returns fixed at their 0.5 km level. The baseline uses cluster-bootstrap mode refits; other rows use all retained posterior draws. The companion full table reports undefined-equilibrium and target-infeasibility shares; target-infeasible draws enter paired contrasts at their best feasible allocation.
\end{minipage}
